\documentclass[11pt]{article}
\usepackage[a4paper,margin=1in]{geometry}
\usepackage{amsmath,amssymb,amsthm}
\usepackage{booktabs}

\newtheorem{theorem}{Theorem}
\newtheorem{lemma}{Lemma}
\newtheorem{proposition}{Proposition}
\newtheorem{corollary}{Corollary}
\newtheorem{conjecture}{Conjecture}
\theoremstyle{definition}
\newtheorem{remark}{Remark}

\newcommand{\Zp}{\mathbb{Z}_p}
\newcommand{\E}{\mathbb{E}}
\newcommand{\Prob}{\Pr}
\newcommand{\cov}{\operatorname{cov}}
\newcommand{\Var}{\operatorname{Var}}

\title{Covering $\Zp$ by Translates of a Random Set of Size $\sqrt p$:\\
The Problem GREEN-081 of Ben Green}
\author{}
\date{August 26, 2026}

\begin{document}
\maketitle

\begin{abstract}
Let $p$ be prime, let $A\subset\Zp$ be a uniformly random subset with
$|A|=k=\lfloor\sqrt p\rfloor$, and let
$\cov(A)=\min\{|B| : A+B=\Zp\}$. GREEN-081 asks whether, almost surely,
$\Zp$ can be covered by $100\sqrt p$ translates of $A$. We separate two
readings. If the translates are themselves independent uniformly random
shifts, we prove the statement asymptotically \emph{false}: writing $X$ for
the number of uncovered points after $m=C\sqrt p$ shifts, we establish the
exact second moment
$\E X^2=pq^m+p(p-1)g_m$ with
$g_m=\E_A\bigl(1-(2k-c_d(A))/p\bigr)^m=q^{2m}(1+O_C(p^{-1}))$, whence
$\Var(X)\le \mu+K(C)\mu^2/p$ for $\mu=\E X=pe^{-C}(1+o(1))$; thus
$X/\mu\to1$ in probability and $(1-o(1))pe^{-100}$ points remain uncovered
with high probability. The sampling threshold lies at
$m^\ast=\sqrt p\,\ln p\,(1+o(1))$, with crossover scale $e^{100}\approx
2.69\times10^{43}$, so the failure, while asymptotically certain, occurs far
beyond computationally accessible primes. Under the intended reading, in
which the translates may be chosen after inspecting $A$, the problem remains
open --- Green records that it is open even with $100$ replaced by $1.01$ ---
and we prove: $\lceil p/k\rceil\le\cov(A)\le(\tfrac12+o(1))\sqrt p\ln p$
deterministically, the upper bound via Lov\'asz--Stein applied to the
$k$-regular translate system, together with
$\cov(A)\le(1+\varepsilon)\sqrt p\ln p$ almost surely. We further prove a
reduction showing that an almost-sure cover follows from an almost-sure
partial cover leaving only $O(\sqrt p)$ points uncovered, and we prove that
Cauchy--Schwarz certificates built solely from first and second moments of
representation functions saturate at a fraction $c/(1+c)<1$ of the group,
which locates the precise obstruction: adaptive control of the tail event
$\{\exists x:\ r_{A+B}(x)=0\}$ past the coupon-collector endgame. Numerical
experiments for primes up to $20011$ confirm the location of the random-shift
threshold, support a Gumbel window law at the critical scale, and place the
greedy benchmark at approximately $0.29\sqrt p\ln p$ chosen translates.
\end{abstract}

\section{The problem and its two readings}\label{sec:model}

Throughout, $p$ denotes a prime, $G=\Zp$, and $k=\lfloor\sqrt p\rfloor\ge2$.
The set $A\subset G$ is drawn uniformly at random from the
$\binom pk$ subsets of cardinality $k$. For $t\in G$ write
$A+t=\{a+t: a\in A\}$. The question

\begin{center}
\emph{``If $A\subset\mathbb Z/p\mathbb Z$ is random with $|A|=\sqrt p$, can we
almost surely cover $\mathbb Z/p\mathbb Z$ with $100\sqrt p$ translates of
$A$?''}
\end{center}

\noindent admits two readings, which we treat separately and which have
different answers:

\begin{itemize}
\item[\textbf{(RA)}] the $100\sqrt p$ translates are themselves independent
uniformly random shifts $t_1,t_2,\dots,t_m$ of $G$, with $m=C\sqrt p$;
\item[\textbf{(RB)}] the translates are selected with knowledge of $A$. Here
the natural object is the covering number
\[
\cov(A):=\min\{|B| :\ B\subset G,\ A+B=G\},
\]
and ``almost surely'' refers to the law of $A$ alone.
\end{itemize}

Green's remark that the problem resists him even when the constant $100$ is
replaced by $1.01$ identifies (RB) as the intended question; reading (RA),
which admits a complete treatment, is analyzed first because it calibrates
what the constant $100$ can and cannot mean. Statements about ``almost sure''
events are asymptotic statements as $p\to\infty$ along primes.

\section{Random translates: exact moments and a threshold theorem}\label{sec:RA}

Fix $m=m(p)\le C_0\sqrt p$, where $C_0$ is an absolute constant, and let
$t_1,\dots,t_m$ be independent uniform points of $G$, independent of $A$.
Write
\[
X:=\Bigl|\,G\setminus\textstyle\bigcup_{i\le m}(A+t_i)\Bigr|,
\qquad q:=1-\frac kp,\qquad \mu:=\E X .
\]

\begin{lemma}[First moment]\label{lem:first}
$\E X=pq^m$ exactly.
\end{lemma}

\begin{proof}
Fix $x\in G$. For a single shift $t$,
$x\in A+t\iff x-t\in A$, and since $t$ is uniform and independent of $A$ the
point $x-t$ is uniform on $G$; hence $\Prob[x\notin A+t]=q$. For distinct
$i,j$ the events $\{x\notin A+t_i\}$, $\{x\notin A+t_j\}$ involve the
independent variables $t_i,t_j$, so
$\Prob[x\text{ uncovered}]=q^m$. Summing over the $p$ choices of $x$ and using
linearity of expectation gives the claim.
\end{proof}

\begin{lemma}[Second moment]\label{lem:second}
For $d=y-x\ne0$ put $c_d(A):=|A\cap(A+d)|$. Then
$(x-A)\cap(y-A)=x-\bigl(A\cap(A+d)\bigr)$, hence
$|(x-A)\cup(y-A)|=2k-c_d(A)$, and conditionally on $A$ a single shift misses
both $x$ and $y$ with probability
\[
w_A(d)=1-\frac{2k-c_d(A)}{p}.
\]
Since $c_d(A+t_0)=c_d(A)$ for every $t_0$, the law of $c_d(A)$ under the
uniform measure on $k$-subsets does not depend on $d\ne0$; writing
\[
g_m:=\E_A\bigl[\,w_A(d)^m\,\bigr],
\]
one has exactly
\[
\E X^2 \;=\; pq^m + p(p-1)\,g_m .
\]
\end{lemma}

\begin{proof}
The diagonal contribution $x=y$ is $p\cdot pq^m/p=pq^m$ by Lemma~\ref{lem:first}.
Fix $x\ne y$ and condition on $A$. By the union--intersection identity above,
a shift $t$ avoids both $x-A$ and $y-A$ precisely when
$t\notin(x-A)\cup(y-A)$, an event of conditional probability $w_A(d)$.
Conditional independence across the $m$ shifts gives
$\Prob[x,y\text{ both uncovered}\mid A]=w_A(d)^m$; averaging over $A$ and
summing over the $p(p-1)$ ordered pairs $(x,y)$ yields the formula. The
$d$-independence of the distribution of $c_d(A)$ follows from the pointwise
identity $c_d(A)=c_d(A+t_0)$ together with invariance of the uniform measure
under translations.
\end{proof}

\begin{remark}[On the average $g_m$ versus its Jensen lower surrogate]\label{rem:jensen}
By concavity of $u\mapsto u^{1/m}$, or equivalently by Jensen's inequality,
$g_m=\E_A[w_A^m]\ge\bigl(\E_A[w_A]\bigr)^m=f^m$, where
\[
f=\frac{\binom{p-2}{k}}{\binom pk}=\frac{(p-k)(p-k-1)}{p(p-1)},
\qquad\quad
\frac{f}{q^{2}}=\frac{p(p-k-1)}{(p-1)(p-k)}
      =1-\frac{k}{(p-k)(p-1)}<1 .
\]
Thus $f^m<q^{2m}$, and indeed $f^m\le g_m$ with strict inequality whenever
$w_A$ is non-constant on the support of $A$ (which holds for $2\le k\le p-2$):
avoiding the pair $\{x,y\}$ is negatively correlated at the level of a single
shift, yet the $m$-th power interacts with the fluctuation of $w_A$ over $A$.
Any computation of $\E X^2$ must therefore average $w_A^m$ over $A$ exactly;
replacing $g_m$ by the surrogate $f^m$ understates the second moment.
\end{remark}

\begin{lemma}[$g_m$ to leading order]\label{lem:gm}
There is a constant $K(C_0)$ such that, uniformly in $m\le C_0\sqrt p$,
\[
g_m=q^{2m}\bigl(1+\theta\bigr),\qquad |\theta|\le \frac{K(C_0)}{p}.
\]
\end{lemma}

\begin{proof}
Exactly,
\[
w_A-q^{2}=\frac{c_d(A)-k^{2}/p}{p},
\]
so with $\Delta_A:=w_A/q^{2}-1$ and $q\ge\tfrac12$ (valid for $p\ge4$),
\[
|\Delta_A|\le\frac{k+1}{pq^{2}}\le\frac{4(k+1)}{p}\le\frac8{\sqrt p}.
\]
Expand $(1+\Delta)^m=1+m\Delta+\rho$ with
$|\rho|\le\tfrac12(m|\Delta|)^{2}(1+|\Delta|)^{m}
\le\tfrac12(m\Delta)^{2}e^{8C_0}$.
First, $\E[c_d]=k(k-1)/(p-1)$, hence
\[
\E\Bigl[c_d-\frac{k^{2}}p\Bigr]=\frac{k(k-1)}{p-1}-\frac{k^{2}}p
=-\frac{k(p-k)}{p(p-1)},
\]
so $m\,|\E\Delta_A|=O(C_0/p)$. Second, $c_d^2=c_d+
\sum_{a\ne b}\mathbf 1_{\{a,a+d,b,b+d\in A\}}$, and the double sum has
expectation at most
$p(p-3)\frac{(k)_4}{(p)_4}+4p\frac{(k)_3}{(p)_3}\le6$, so $\E[c_d^2]\le C_1$
absolutely and $m^{2}\E[\Delta_A^{2}]=O(C_0^{2}/p)$. Taking expectations in
the expansion combines the two estimates.
\end{proof}

\begin{theorem}[Variance bound]\label{thm:var}
Uniformly in $m\le C_0\sqrt p$,
\[
\Var(X)=\mu\bigl(1-q^{m}\bigr)+p(p-1)\bigl(g_m-q^{2m}\bigr)
\;\le\;\mu+K'(C_0)\,\frac{\mu^{2}}{p}.
\]
Consequently, if $\mu\to\infty$ then $X/\mu\to1$ in $L^{2}$ and
$\Prob[X=0]\le\Var(X)/\mu^{2}\to0$.
\end{theorem}

\begin{proof}
Substitute Lemma~\ref{lem:gm} into Lemma~\ref{lem:second} and subtract
$(\E X)^2=p^{2}q^{2m}$; the sign of $g_m-q^{2m}$ is immaterial since
$p(p-1)|g_m-q^{2m}|\le K(C_0)\,p^{-1}p^{2}q^{2m}=K'(C_0)\mu^{2}/p$.
The concentration statement is Chebyshev's inequality applied to
$\{X=0\}\subseteq\{|X-\mu|\ge\mu\}$.
\end{proof}

\begin{theorem}[Threshold theorem; the random-translates reading is false]\label{thm:threshold}
Let $m\le C_0\sqrt p$.
\begin{enumerate}
\item[(i)] Fix any constant $C>0$ and set $m=C\lfloor\sqrt p\rfloor$. Then
\[
\mu=p\Bigl(1-\frac1{\sqrt p}\Bigr)^{C\sqrt p}=pe^{-C}(1+o(1))\longrightarrow\infty ,
\]
so w.h.p.\ $X\ge(1-o(1))\,pe^{-C}$ and
$\Prob[X=0]\to0$. In particular, for $C=100$, random $100\sqrt p$ translates
fail to cover $\Zp$ with probability tending to one: under reading (RA) the
statement of GREEN-081 is false.
\item[(ii)] If instead $m\ge\sqrt p\,(\ln p+\omega(1))$ for some
$\omega(1)\to\infty$, then by
$\ln(1-\tfrac1{\sqrt p})=-\tfrac1{\sqrt p}-\tfrac1{2p}+O(p^{-3/2})$,
\[
\mu\le p\exp\bigl(-(\ln p+\omega(1))(1-o(1))\bigr)=e^{-\omega(1)+o(1)}\to0 ,
\]
and Markov's inequality gives $\Prob[\text{full cover}]\to1$.
\item[(iii)] Outside the critical window
$|m/\sqrt p-\ln p-c|=O(1)$ (see Conjecture~\ref{conj:gumbel}), random shifts
cover w.h.p.\ if and only if $m/\sqrt p-\ln p\to+\infty$ and fail w.h.p.\ if
$m/\sqrt p-\ln p\to-\infty$: the threshold sits at
$m^{\ast}\approx\sqrt p\ln p$, not at $C\sqrt p$ for any fixed constant $C$.
\end{enumerate}
\end{theorem}

\begin{corollary}[Crossover scale]\label{cor:crossover}
Under reading (RA), covers succeed w.h.p.\ for $p\le e^{100-\omega(1)}$ and
fail w.h.p.\ for $p\ge e^{100+\omega(1)}$, where
$e^{100}=2.688\times10^{43}$. At the accessible primes $p=997$ and $p=10007$,
evaluating the exact formula of Lemma~\ref{lem:first} at
$m=\lfloor100\sqrt p\rfloor$ gives
$\mu=7.246049\times10^{-42}$ resp.\ $2.262069\times10^{-40}$: the asymptotic
failure is invisible to direct computation at feasible sizes, while the
threshold of Theorem~\ref{thm:threshold}(iii) lies far below
$100\sqrt p$ and is fully observable (Section~\ref{sec:numerics}).
\end{corollary}

\begin{conjecture}[Gumbel window]\label{conj:gumbel}
If $m/\sqrt p-\ln p\to c\in\mathbb R$, then
$X$ converges in distribution to $\mathrm{Poisson}(e^{-c})$ and
$\Prob[\text{full cover}]\to\exp(-e^{-c})$.
\end{conjecture}

\begin{remark}[Status of the conjecture]\label{rem:window}
For each fixed $j$ the factorial moments satisfy
$\E[(X)_j]=(p)_j h_j$ where $h_j=\E_A[u_j(A)^m]$ and $u_j(A)$ is the
probability that a single shift misses $j$ prescribed points; the expansion
of Lemma~\ref{lem:gm} extends to $h_j=(1-jk/p)^{m}(1+o(1))$ for fixed $j$,
using boundedness of joint autocorrelation moments of order at most $2j$.
Together with $g_m/q^{2m}\to1$ this makes $\mathrm{Poisson}(e^{-c})$ the
consistent limit heuristic. The obstruction to a proof is that the events
$\{x\text{ uncovered}\}$, $x\in G$, share \emph{all} of the randomness, so the
dependency graph is complete and Chen--Stein bounds of Arratia--Goldstein--Gordon
type are vacuous ($b_2\simeq\lambda^2\not\to0$); a genuinely finer Stein
argument would be required.
\end{remark}

\section{Chosen translates: what is provable, and where it breaks}\label{sec:RB}

We now turn to the intended reading. Recall $\cov(A)=\min\{|B|:A+B=G\}$.

\begin{proposition}[Counting lower bound]\label{prop:lower}
$\cov(A)\ge\lceil p/k\rceil$ for every $A\subset G$ with $|A|=k$. Moreover no
exact factorization $G=A\oplus B$ with $|A|,|B|\ge2$ exists, by the
Haj\'os--de Bruijn theorem on factorizations of cyclic groups of prime order;
a perfect tiling is impossible, so any near-optimal cover must tolerate
controlled overlap.
\end{proposition}

\begin{proof}
$p=|A+B|\le|A||B|=k|B|$.
\end{proof}

\begin{proposition}[Almost-sure upper bound]\label{prop:as}
For every fixed $\varepsilon>0$,
\[
\Prob_A\Bigl[\cov(A)\le(1+\varepsilon)\sqrt p\,\ln p\Bigr]\longrightarrow1 .
\]
\end{proposition}

\begin{proof}
Put $M=\lceil(1+\varepsilon)\sqrt p\ln p\rceil$ and draw $B$ uniformly from
the $M$-subsets of $G$, independently of $A$. Conditionally on $A$, the
expected number of uncovered points is
$p\cdot\Prob[x\notin A+B]$, and for fixed $x$ the probability that the uniform
$M$-subset $x-B$ avoids the $k$-set $A$ is at most $(1-k/p)^M$, so
\[
\E\bigl[\#\text{uncovered}\bigr]
=p\,\frac{\binom{p-M}{k}}{\binom pk}
\le p\Bigl(1-\frac kp\Bigr)^{M}
\le p\,e^{-M/\sqrt p}
=p^{-(\varepsilon-o(1))}\longrightarrow0 .
\]
If $\cov(A)>M$ then \emph{every} $M$-subset fails to cover, in particular the
random one, so pointwise in $A$,
$\mathbf 1_{\{\cov(A)>M\}}\le\Prob_B[\text{fail}\mid A]$. Averaging over $A$
yields the claim.
\end{proof}

\begin{proposition}[Deterministic Lov\'asz--Stein upper bound]\label{prop:lovasz}
Every $A\subset G$ with $|A|=k$ satisfies
\[
\cov(A)\le\frac pk\,(1+\ln k)=\Bigl(\tfrac12+o(1)\Bigr)\sqrt p\,\ln p .
\]
Note the deterministic constant $\tfrac12$ improves on the probabilistic
constant $1$ of Proposition~\ref{prop:as}.
\end{proposition}

\begin{proof}[Proof sketch]
Consider the set system $\{A+t\}_{t\in G}$ on the ground set $G$. It is
$k$-uniform and $k$-regular: each $x\in G$ lies in exactly the $k$ translates
$A+(x-a)$, $a\in A$. Regularity forces the fractional covering number to equal
$\tau^\ast=p/k$, and Lov\'asz's theorem
$\tau\le\tau^\ast(1+\ln d)$ for $d$-uniform systems (see Lov\'asz 1975;
S.\,K.\,Stein 1974; Chv\'atal 1979) gives the bound.
\end{proof}

\begin{proposition}[Ceiling for mean-square methods]\label{prop:energy}
Let $r=r_{A+B}$ denote the representation function and define the
Cauchy--Schwarz certificate
$\mathcal C(A,B):=\bigl(\sum_x r(x)\bigr)^{2}\big/\sum_x r(x)^{2}\le|A+B|$.
For $|B|=c\sqrt p$,
\[
\E\sum_x r(x)=cp,
\qquad
\E\sum_xr(x)^{2}=cp+c^{2}p+o(p),
\qquad\text{hence}\qquad
\mathcal C(A,B)\le\frac{cp}{1+c}<p .
\]
Moreover $\E\bigl[\mathsf E(A)\bigr]=\E\sum_dc_d(A)^{2}=2p+o(p)$, essentially
minimal for a $k$-set since always
$\mathsf E(A)\ge\max\bigl(k^{2},k^{4}/|A+A|\bigr)=p+o(p)$, and this minimal
energy carries no information about the zero set
$\{x:r(x)=0\}$. Consequently, no argument that consumes only first and second
moments of representation counts can certify full coverage at
$M=O(\sqrt p)$ translates: such arguments saturate at fraction
$c/(1+c)<1$ of the group.
\end{proposition}

\begin{proof}
Deterministically $\sum_xr(x)=k|B|=cp$. For the second moment,
$\sum_xr(x)^{2}=\sum_xr(x)+\sum_xr(x)(r(x)-1)$, and the second term counts
ordered pairs of distinct translates meeting at $x$:
\[
\E\sum_xr(x)\bigl(r(x)-1\bigr)
=M(M-1)\,\E[c_d]
=M(M-1)\frac{k(k-1)}{p-1}=c^{2}p+o(p),
\]
using $\E[c_d]=k(k-1)/(p-1)$. For the energy,
$c_d^{2}=c_d+\sum_{a\ne b}\mathbf 1_{\{a,a+d,b,b+d\in A\}}$ with the double
sum of expectation $O(1)$ per $d$ (as in the proof of Lemma~\ref{lem:gm}),
whence $\E\mathsf E(A)=\sum_d\E c_d^{2}=p\cdot(2+o(1))$. The lower bounds
$\mathsf E(A)\ge k^{2}$ (diagonal) and
$\mathsf E(A)\ge k^{4}/|A+A|\ge k^{4}/p$ (Cauchy--Schwarz) give minimality.
Finally $\mathcal C$ is Cauchy--Schwarz applied to the indicator of the
support of $r$ against $r$ itself.
\end{proof}

\begin{theorem}[Completion lemma and reduction]\label{thm:reduction}
Let $B_0\subset G$ be any partial cover with leftover set $U$,
$L:=|U|=p-|A+B_0|$. Then:
\begin{enumerate}
\item[(a)] $\cov(A)\le |B_0|+L$; indeed each $x\in U$ is covered by the single
translate $A+(x-a_0)$ for any fixed $a_0\in A$.
\item[(b)] Greedy completion costs at most
$\bigl\lceil\frac pk\ln(L+1)\bigr\rceil+1$ further translates: at leftover
level $U'$, $\sum_{t}|(A+t)\cap U'|=k|U'|$ exactly, so some translate adds at
least $k|U'|/p$ fresh points and $|U'|$ decays geometrically by the factor
$1-k/p\le e^{-k/p}$ per step.
\end{enumerate}
Consequently $\cov(A)=O(\sqrt p)$ w.h.p.\ as soon as w.h.p.\ there exist
$O(\sqrt p)$ translates covering all but $O(\sqrt p)$ points of $G$; and
GREEN-081(RB) is equivalent to the assertion
\[
\inf_{B\subset G}\bigl(|B|+L_B\bigr)=O(\sqrt p)\ \text{w.h.p.},\qquad
L_B:=p-|A+B|.
\]
Within the completion family (b) the total equals
$M+\sqrt p\ln(L+1)(1+o(1))$, which is $O(\sqrt p)$ only if
$\ln L=O(1)$: the logarithmic endgame --- eliminating the last few uncovered
points --- is exactly where every available method pays $\sqrt p\ln p$ rather
than $\sqrt p$.
\end{theorem}

\subsection*{Structural obstructions}

\paragraph{Disjoint-packing (design) routes.}
Hypergraph matching theory --- the R\"odl nibble, Pippenger--Spencer, Keevash's
existence of designs --- produces \emph{disjoint} edge families. In the present
translate system two edges intersect in $c_{t-t'}(A)$ points with
$\E[c_d]=k(k-1)/(p-1)=1+O(p^{-1/2})$: typical translates already meet in
$\Theta(1)$ points, disjoint edge families are structurally absent, and no
general ``bounded-overlap approximate decomposition with quantitative
leftover'' theorem is available for Cayley hypergraphs. (Standard heuristics
give $|A-A|=(1-e^{-1}+o(1))p$ w.h.p.; the per-shift avoidance probability
$\Prob[d\notin A-A]$ equals the probability that a uniform $k$-set is
independent in the difference cycle of step $d$, namely
$\frac{p}{p-k}\binom{p-k}{k}/\binom pk\to e^{-1}$.)

\paragraph{No combinatorial obstruction to efficiency.}
If $|B|=\beta\sqrt p$ and multiplicities $r(x)\ge1$ held everywhere, flat
multiplicity would force total pairwise collisions
$\sum_x\binom{r(x)}{2}\gtrsim\frac{\beta^{2}-\beta}{2}p$, whereas typical
translates supply collision budget
$\sum_{t<t'}|(A+t)\cap(A+t')|
\approx\binom{\beta\sqrt p}{2}\frac{k^{2}}{p}
=\frac{\beta^{2}}{2}p+o(p)$, and
$\frac{\beta^{2}}2p\ge\frac{\beta^{2}-\beta}2p$ always. Nothing structurally
forbids an efficient near-tiling; the difficulty is entirely in the tail
behavior of adaptively chosen families.

\paragraph{The precise gap.}
What a proof must deliver is control of
$\Prob[\exists x:\ r_{A+B}(x)=0]$ for an explicit adaptive $B$ with
$|B|=O(\sqrt p)$. A union bound loses $pe^{-c}$; under \emph{random} $B$ the
bad events behave nearly independently --- that is precisely the mechanism of
Conjecture~\ref{conj:gumbel} --- so adaptivity must suppress the joint tail.
No current technique (second moments, Janson inequalities, Chen--Stein theory,
entropy methods, containers, LP rounding) provides this in a
translation-invariant setting, and we are aware of no reduction of
the statement to another named conjecture beyond Green's original formulation.

\section{Numerical experiments}\label{sec:numerics}

All computations use exact integer coverage checks (boolean incidence arrays
for unions; integer-rounded convolutions, validated against brute-force
recounts, for scores). Randomness comes from fixed documented seeds, so every
table reproduces exactly. Primes examined:
$p\in\{101,199,401,599,797,997,1999,4003,5003,10007,20011\}$.

\subsection{Location of the random-shift threshold}
For each prime, many independent trials were run; in each trial a fresh $A$
was drawn and shifts were added one at a time until $G$ was covered exactly,
recording the stopping time $m^\ast$ (trials capped at
$\lceil\sqrt p(\ln p+7)\rceil$, an event of negligible probability).

\begin{center}
\begin{tabular}{rrrrr}
\toprule
$p$ & $k$ & mean $m^\ast$ & median & mean$/\sqrt p\ln p$\\
\midrule
101 & 10 & 49.71 & 48 & 1.0717\\
199 & 14 & 80.64 & 77 & 1.0799\\
401 & 20 & 128.31 & 124 & 1.0690\\
599 & 24 & 170.93 & 165 & 1.0920\\
797 & 28 & 203.97 & 199 & 1.0814\\
997 & 31 & 237.67 & 231 & 1.0901\\
1999 & 44 & 374.96 & 356 & 1.1034\\
4003 & 63 & 563.91 & 542 & 1.0745\\
5003 & 70 & 631.75 & 617 & 1.0486\\
10007 & 100 & 964.62 & 937 & 1.0469\\
20011 & 141 & 1435.62 & 1403 & 1.0247\\
\bottomrule
\end{tabular}
\end{center}

The ratios drift toward $1$ at the rate $1+\gamma/\ln p$
($1.125$ at $p=101$, $1.058$ at $p=20011$), matching
Theorem~\ref{thm:threshold}(iii): the threshold is $\sqrt p\ln p$, while
$100\sqrt p$ exceeds it tenfold at $p=20011$.

\subsection{Behavior at the critical window}
At budgets $m=\lceil\sqrt p(\ln p+c)\rceil$, the empirical success frequency
compares with the Gumbel prediction $\exp(-e^{-c})$:

\begin{center}
\begin{tabular}{rrrrr}
\toprule
$c$ & predicted & $p=997$ & $p=1999$ & $p=5003$\\
\midrule
$-1$ & 0.0660 & 0.0625 & 0.0667 & 0.0933\\
$0$ & 0.3679 & 0.3600 & 0.3667 & 0.4467\\
$1$ & 0.6922 & 0.6900 & 0.6800 & 0.7400\\
$2$ & 0.8734 & 0.8650 & 0.7867 & 0.9000\\
\bottomrule
\end{tabular}
\end{center}

Counting uncovered points after exactly $m=\mathrm{round}(\sqrt p(\ln p+c))$
shifts (300 trials per cell), the empirical mean tracks $e^{-c}$ (at $p=997$:
$0.977,0.390,0.150$ for $c=0,1,2$ against predictions $1,0.368,0.135$),
variance tracks the mean, and the frequency of zero uncovered points brackets
$\exp(-e^{-c})$ throughout --- all consistent with
Conjecture~\ref{conj:gumbel}.

\subsection{The regime of the literal constant}
At $m=\lfloor100\sqrt p\rfloor$ for $p\in\{997,10007\}$, no uncovered point
was ever observed, in accord with Corollary~\ref{cor:crossover}:
the theoretical means $7.246\times10^{-42}$ and
$2.262\times10^{-40}$ preclude observation at feasible scales.

\subsection{Chosen translates: greedy benchmark}
A greedy experiment emulates reading (RB): starting from $u=G$ uncovered,
repeatedly select a translate maximizing $|(A+t)\cap u|$ until $u=\emptyset$,
recording the step count $g$. Scores were computed exactly (integer-rounded
convolutions cross-checked against direct recounts).

\begin{center}
\begin{tabular}{rrrrrr}
\toprule
$p$ & trials & $g$ mean & $g/(\sqrt p\ln p)$ &
$g/(\tfrac12\sqrt p\ln p)$ & Lov\'asz--Stein bound\\
\midrule
101 & 200 & 15.55 & 0.3353 & 0.671 & 33.4\\
997 & 60 & 65.42 & 0.3000 & 0.600 & 142.6\\
1999 & 40 & 100.13 & 0.2946 & 0.589 & 217.4\\
5003 & 25 & 175.20 & 0.2908 & 0.582 & 375.1\\
20011 & 8 & 402.88 & 0.2876 & 0.575 & 844.3\\
\bottomrule
\end{tabular}
\end{center}

Greedy outperforms the proven ceiling of Proposition~\ref{prop:lovasz} by a
factor ${\approx}2.1$, yet $g/\sqrt p$ grows like $0.29\ln p$: even the best
heuristic visible to local search carries the logarithm, and the endgame adds
only $O(1)$ points per step ($\approx2$ in the largest run) --- exactly the
saturation predicted by the discussion surrounding
Theorem~\ref{thm:reduction}.

\section{Verdict}\label{sec:verdict}

\begin{enumerate}
\item \textbf{Reading (RA), random translates: the statement is disproven,
asymptotically.} By Theorem~\ref{thm:var} and
Theorem~\ref{thm:threshold}(i), random $100\sqrt p$ translates leave
$(1-o(1))pe^{-100}$ points uncovered with probability tending to one; covers
occur w.h.p.\ only for $p\lesssim e^{100}$
(Corollary~\ref{cor:crossover}). The proof is elementary and
self-contained: exact first and second moments
(Lemmas~\ref{lem:first}--\ref{lem:second}, with the exact averaged quantity
$g_m=\E_A[w_A^m]$ of Lemma~\ref{lem:gm}), the variance bound, and Chebyshev.
\item \textbf{Reading (RB), chosen translates: open.} Established here:
\[
\Bigl\lceil\frac pk\Bigr\rceil\;\le\;\cov(A)\;\le\;
\Bigl(\tfrac12+o(1)\Bigr)\sqrt p\ln p \quad\text{for every }A,
\qquad
\cov(A)\le(1+\varepsilon)\sqrt p\ln p\ \text{a.s.},
\]
together with the reduction of Theorem~\ref{thm:reduction}. The exact missing
bridge is
\[
\Prob_A\bigl[\exists B:\ |B|=O(\sqrt p),\ L_B=O(\sqrt p)\bigr]\longrightarrow1 ;
\]
by Proposition~\ref{prop:energy} no mean-square or additive-energy argument
can certify this (certificates cap at fraction $c/(1+c)$ of the group), by the
structural discussion in Section~\ref{sec:RB} disjoint-packing/design
machinery is inapplicable, and every completion analysis pays the endgame
logarithm. Numerically, the greedy benchmark sits at
$\approx0.29\sqrt p\ln p$: between the proven ceiling and the conjectured
$O(\sqrt p)$.
\item Numerics confirm the threshold location
(Theorem~\ref{thm:threshold}(iii)), support the Gumbel window law
(Conjecture~\ref{conj:gumbel}), confirm the practical invisibility of the
(RA) failure below $p\sim e^{100}$, and quantify the greedy benchmark.
\end{enumerate}

\section*{References}
\begin{itemize}
\item B. Green, \emph{100 Open Problems in Additive Combinatorics}, Problem 39.
\item L. Lov\'asz, On the ratio of optimal integral and fractional covers,
\emph{Discrete Math.} \textbf{13} (1975), 383--390.
\item S. K. Stein, Two combinatorial covering theorems, \emph{L'Enseignement
Math\'ematique} \textbf{20} (1974).
\item V. Chv\'atal, A greedy heuristic for the set-cover problem,
\emph{Math. Oper. Res.} \textbf{4} (1979), 233--235.
\item G. Haj\'os, \"Uber einfache und mehrfache Bedeckung von
$n$-dimensionalen R\"aumen mit einem W\"urfelgitter, \emph{Math. Z.}
\textbf{47} (1942), 427--467.
\item N. G. de Bruijn, On the factorization of finite abelian groups,
\emph{Indag. Math.} \textbf{16} (1953), 254--258.
\end{itemize}

\end{document}
